% Source: https://artofproblemsolving.com/community/c3h2658474_the_problem_of_numbers
% Stable identity: aops-topic-2658474-q01
Group the elements of $S$ by their number $d\ge1$ of base-$n$ digits.
\begin{itemize}
\item Each of the $d$ digits avoids $k$, so there are at most $(n-1)^d$ elements of $S$ with $d$ digits.
\item Every $d$-digit number is at least $n^{d-1}$, so each such reciprocal is at most $\dfrac1{n^{d-1}}$.
\end{itemize}
Hence the $d$-digit elements contribute at most $\dfrac{(n-1)^d}{n^{d-1}}$, and since all terms are positive,
\[\sum_{s\in S}\dfrac1s\le\sum_{d=1}^{\infty}\dfrac{(n-1)^d}{n^{d-1}}=(n-1)\sum_{d=0}^{\infty}\left(\dfrac{n-1}{n}\right)^d=(n-1)\cdot\dfrac{1}{1-\dfrac{n-1}{n}}=n(n-1).\]
The geometric series converges because $0<\dfrac{n-1}{n}<1$. So the partial sums of $\sum_{s\in S}\dfrac1s$ are bounded, the series converges, and its sum is $\boxed{\le n(n-1)}$. $\blacksquare$

(For example, the integers with no digit $9$ in base ten have reciprocal sum at most $90$; its actual value is about $22.92$. No simple closed form for these sums is known.)
